% §8.1 Shannon Entropy Refresher
%
% Cross-reference resolution status:
%
% DEFINED here:
%   \label{sec:entropy-refresher}                 -- this section
%
% ALREADY DEFINED (chapter 2):
%   \Cref{sec:bsc}                                -- chapter 2 §2.1
%
% FORWARD REFS within chapter 8 (resolve in later Plan 8 tasks):
%   \Cref{sec:bsc-capacity}                       -- §8.2 (Task 3)
%   \Cref{sec:entropy-floor}                      -- §8.3 (Task 4)
%   \Cref{sec:bloom-gap}                          -- §8.4 (Task 5)
%
% §8.1 collects the three pieces of Shannon's information theory used
% by the rest of chapter 8: the entropy $H(X)$, the binary entropy
% function $H(p)$, and the mutual information $I(X; Y)$. Stated, not
% derived; Cover and Thomas is the depth reference.

\section{Shannon Entropy Refresher}
\label{sec:entropy-refresher}

The chapter's two headline results, the BSC capacity in
\Cref{sec:bsc-capacity} and the entropy floor for approximate sets in
\Cref{sec:entropy-floor}, both rest on a small slice of Shannon's
information theory: the entropy of a random variable, the special
form of that entropy for a single binary outcome, and the mutual
information between two random variables. This section collects the
three pieces and fixes notation. Cover and Thomas~\cite{cover-thomas}
is the standard depth reference; nothing here is derived from first
principles.

\subsection*{Entropy of a discrete random variable}

For a discrete random variable $X$ taking values in a finite alphabet
$\mathcal{X}$ with distribution $\Pr[X = x]$, the \emph{Shannon
entropy} of $X$ is
\begin{equation}
\label{eq:entropy}
  H(X) \;=\; -\sum_{x \in \mathcal{X}} \Pr[X = x] \, \log_2 \Pr[X = x],
\end{equation}
with the convention $0 \log_2 0 = 0$. The base-$2$ logarithm sets the
unit to \emph{bits}. The operational reading is that $H(X)$ is the
average number of yes/no questions an optimal binary search needs to
identify a draw of $X$: at one extreme, a constant random variable
has $H(X) = 0$ (no questions needed); at the other extreme, a uniform
distribution over $|\mathcal{X}|$ outcomes has $H(X) = \log_2
|\mathcal{X}|$ bits (a balanced binary search). Real-world
distributions sit somewhere between. The quantity $H(X)$ is also the
tight lower bound on the expected length of any prefix-free binary
code for $X$, by Shannon's source-coding theorem, and that
operational reading is what \Cref{sec:entropy-floor} uses.

\subsection*{The binary entropy function}

The single most-used special case is a Bernoulli random variable.
For $X \sim \mathrm{Bernoulli}(p)$ with $\Pr[X = 1] = p$ and
$\Pr[X = 0] = 1 - p$, equation~\eqref{eq:entropy} collapses to a
function of $p$ alone:
\begin{equation}
\label{eq:binary-entropy}
  H(p) \;=\; -p \log_2 p \;-\; (1 - p) \log_2 (1 - p),
  \qquad p \in [0, 1].
\end{equation}
We overload the symbol $H$: the form $H(X)$ takes a random variable
and the form $H(p)$ takes a scalar $p \in [0, 1]$; context
disambiguates. The function $H(p)$ has three properties used
repeatedly in this chapter and in chapter~14:

\begin{enumerate}
  \item \emph{Concave on $[0, 1]$.} The second derivative
    $H''(p) = -1 / (p (1 - p) \ln 2)$ is negative on $(0, 1)$, so $H$
    is strictly concave. Jensen-style averaging arguments apply.
  \item \emph{Maximum $1$ at $p = 1/2$.} Setting $H'(p) = 0$ gives
    $\log_2 ((1 - p) / p) = 0$, so $p = 1/2$, where
    $H(1/2) = 1$ bit. A fair coin is the most uncertain single bit.
  \item \emph{Zero at the boundaries.} $H(0) = H(1) = 0$ by the
    $0 \log_2 0 = 0$ convention. A deterministic bit carries no
    information.
\end{enumerate}

\noindent The function is symmetric about $p = 1/2$, since
$H(p) = H(1 - p)$ swaps the two terms. Numerically, $H(0.01) \approx
0.081$ bit, $H(0.1) \approx 0.469$ bit, $H(0.5) = 1$ bit. The
quantity that matters in \Cref{sec:bsc-capacity} is $H(\fprate)$
where $\fprate$ is the BSC crossover rate; the quantity that matters
in \Cref{sec:entropy-floor} is the related quantity $-\log_2 \fprate$,
which is the entropy of an indicator that an event of probability
$\fprate$ occurs (the dominant term of $H(\fprate)$ as $\fprate \to
0$).

\subsection*{Mutual information}

When two random variables $X$ and $Y$ are jointly distributed, the
\emph{mutual information}
\begin{equation}
\label{eq:mutual-info}
  I(X; Y) \;=\; H(X) - H(X \mid Y) \;=\; H(Y) - H(Y \mid X)
\end{equation}
measures how much observing $Y$ reduces uncertainty about $X$, or
equivalently how much observing $X$ reduces uncertainty about $Y$.
The two forms in \eqref{eq:mutual-info} are equal: $I(X; Y)$ is
symmetric in its two arguments, and the chapter uses both
orientations as convenience dictates. Here $H(X \mid Y) = \sum_y
\Pr[Y = y] \, H(X \mid Y = y)$ is the conditional entropy of $X$
given $Y$, the average residual uncertainty about $X$ after $Y$ is
known. An equivalent form writes the mutual information directly
from the joint distribution,
\begin{equation}
\label{eq:mutual-info-joint}
  I(X; Y) \;=\; \sum_{x, y} \Pr[X = x, Y = y]
    \, \log_2 \frac{\Pr[X = x, Y = y]}{\Pr[X = x] \Pr[Y = y]},
\end{equation}
showing that $I(X; Y)$ is a Kullback-Leibler divergence between the
joint distribution and the product of marginals, and so always
nonnegative. Equality $I(X; Y) = 0$ holds iff $X$ and $Y$ are
independent: the joint factorizes, the ratio inside the logarithm
is uniformly $1$, and observing one of the pair tells you nothing
about the other. At the opposite extreme, if $X$ is a function of
$Y$, then $H(X \mid Y) = 0$ and $I(X; Y) = H(X)$: the observation of
$Y$ pins down $X$ completely, and the reduction equals all of $X$'s
prior uncertainty. Between those two endpoints, $I(X; Y)$
interpolates smoothly and is the central object that
\Cref{sec:bsc-capacity} optimizes.

For a noisy channel with input $X$ and output $Y$ governed by a
fixed conditional distribution $\Pr[Y \mid X]$, the mutual
information $I(X; Y)$ depends on the input distribution
$\Pr[X]$. The \emph{channel capacity} is the maximum:
\begin{equation}
\label{eq:capacity-def}
  C \;=\; \max_{\Pr[X]} I(X; Y).
\end{equation}
This is the most information that a single use of the channel can
convey, optimized over what the sender is allowed to choose (the
distribution on inputs). \Cref{sec:bsc-capacity} works out
\eqref{eq:capacity-def} explicitly for the BSC and arrives at the
compact form $C = 1 - H(\fprate)$.

\subsection*{The chapter's plan for these tools}

Equations~\eqref{eq:entropy}, \eqref{eq:binary-entropy},
\eqref{eq:mutual-info}, and~\eqref{eq:capacity-def} are the only
information-theoretic objects this chapter needs.
\Cref{sec:bsc-capacity} applies them to the BSC of
\Cref{sec:bsc} and derives the capacity $C = 1 - H(\fprate)$, giving
the chapter-2 channel its quantitative information-theoretic summary
and recovering the familiar endpoints (one bit per use of a noiseless
channel, zero bits per use of a maximally noisy channel).
\Cref{sec:entropy-floor} then applies the source-coding reading of
$H(X)$ to the bit-string representation of an approximate set: the
result is a per-element entropy floor of $-\log_2 \fprate$ bits for
any approximate set with zero false-negative rate and false-positive
rate at most $\fprate$. The reading is direct. An approximate-set
representation is a binary code for some random object, and any such
code is bound below by the entropy of what it encodes; the floor is
that entropy made explicit. \Cref{sec:bloom-gap} compares the
standard Bloom filter's actual cost to that floor and recovers the
literature's $\log_2 e \approx 1.443$ overhead constant. Throughout,
the role of \eqref{eq:binary-entropy} is to bridge a probability
into the bit budget that probability implies, and the role of
\eqref{eq:capacity-def} is to bridge a channel into the per-use
information bound it sets.
